\subsection{Geological Media as Physical Materials}

Geological media are natural Earth materials considered as physical bodies that can store, transmit, and respond to mechanical and thermal disturbances. In the context of thermoelastic wave propagation, this term includes both the rock material itself and, at a larger scale, the rock mass, that is, the rock together with its pore space, contained fluids, fractures, joints, and other structural discontinuities. Geomechanics therefore does not treat rocks as chemically simple or structurally uniform solids. Instead, rocks are treated as geomaterials whose behavior reflects a complex combination of solid phases, pore space, possible fluids, microstructure, and geological history. This point is fundamental because the physical response of a geological medium depends not only on the minerals it contains, but also on how those minerals are arranged, connected, fractured, cemented, or weathered in space \cite{jaeger2007,mavko2009,robertson1988}.

Unlike idealized engineering materials manufactured under controlled conditions, natural rocks are products of geological processes such as magmatic crystallization, sediment deposition, burial, compaction, cementation, recrystallization, deformation, uplift, and weathering. As a result, rocks and rock masses are rarely perfectly homogeneous. Even when a rock appears uniform in a hand specimen, it usually consists of several mineral phases with different crystal structures and different elastic and thermal properties. In addition, the rock mass may contain pores, microcracks, healed cracks, open fractures, bedding surfaces, foliation planes, stylolitic seams, or weathered zones that interrupt the continuity of the material and may make its properties direction-dependent. Rock mechanics therefore differs from the classical strength theory of perfectly continuous solids because real rock masses often lie between two end members: a continuous medium at one scale and a discontinuous medium at another \cite{jaeger2007,schultz2019}.

Mineral composition is one of the first-order controls on the physical character of geological media. Rocks are aggregates of mineral grains, and each mineral contributes its own density, stiffness, thermal conductivity, heat capacity, and thermal expansion behavior. Quartz-rich rocks, carbonate rocks, and mafic crystalline rocks do not conduct heat or deform elastically in the same way because quartz, calcite, feldspar, mica, pyroxene, and amphibole differ in their physical properties. Rock texture then adds another level of control. Grain size, grain shape, packing, sorting, cement type, grain contacts, and pore geometry all influence bulk density, effective porosity, stiffness, permeability, and wave velocity. In porous media, not all void space has the same physical meaning: effective porosity refers to the interconnected pore space available for fluid transmission, and the geometry of that connected space may matter as much as the total pore volume itself. Classical rock-physics and hydrologic-property compilations show that porosity is affected by grain-size distribution, grain shape, deposition, compaction, and solidification, while pore fluids and pore-space stiffness strongly influence rock modulus and seismic velocity \cite{wolff1982,mavko2009,robertson1988}.

The same rock body may therefore exhibit substantial spatial variability. A sandstone bed may be loosely cemented in one interval and tightly cemented in another. A basalt flow may contain dense massive interiors, vesicular upper zones, and cooling joints. A granite may be fresh and intact at depth but strongly jointed and weathered near the surface. A limestone may pass from compact micritic layers to fractured and vuggy zones over short distances because of diagenesis and dissolution. Such internal variation means that physical properties measured in the laboratory on a small core sample are not automatically representative of the whole field-scale medium. Rock-property compilations repeatedly emphasize the importance of scale, the representativeness of point measurements, and the dependence of porosity and permeability on both matrix structure and fractures. For the same reason, field geology treats bedding, foliation, fracture sets, and weathering zones not as secondary details, but as structures that can reorganize the pathways of stress, heat, and fluids within the rock mass \cite{wolff1982,robertson1988,schultz2019}.

To analyze such complex media, theoretical descriptions often use simplifying categories. A homogeneous medium is one whose averaged properties do not vary from point to point within the domain of interest, whereas a heterogeneous medium shows spatial variation in composition or properties. An isotropic medium has the same relevant property in all directions, while an anisotropic medium shows directional dependence. A continuous medium is one for which the material can be represented by smoothly varying field quantities such as temperature, displacement, strain, and stress. A discontinuous medium, by contrast, contains explicit internal breaks such as fractures, joints, faults, or interfaces that may need to be treated separately. In geomechanics, these categories are not absolute geological truths but modeling assumptions whose usefulness depends on scale. Rocks can often be approximated as continua only when the observation scale is sufficiently larger than the characteristic size of grains, pores, and small discontinuities. At the same time, large-scale bedding, foliation, or fracture networks may remain too important to average out and must be retained explicitly \cite{cornet2007,jaeger2007}.

These assumptions matter directly for thermoelastic wave propagation. If a geological medium is treated as homogeneous and isotropic, wave speeds and thermal diffusion can be described with averaged density, elastic moduli, thermal conductivity, and heat capacity. Such a formulation is mathematically convenient and is often appropriate as a first approximation. However, when layering, mineral alignment, fracture orientation, stress-induced crack closure, or preferred pore geometry are significant, propagation becomes directional. In that case, the same thermal disturbance or elastic pulse may travel at different speeds, attenuate differently, or induce different stress concentrations depending on direction. Seismic anisotropy studies show that directional behavior is not exceptional in Earth materials; it may arise from aligned minerals, horizontal layering, fracture sets, and stress-induced crack fabrics. Thus, even in an introductory theoretical treatment, the continuum and isotropy assumptions must be understood as useful idealizations rather than complete descriptions of geological reality \cite{aki2002,crampin1981,thomsen1986,schultz2019}.

For this reason, the study of thermoelastic wave propagation in rocks begins with a careful physical-geological view of the medium itself. Before discussing the governing equations of elasticity, heat conduction, or thermoelastic stress, it is necessary to recognize that geological materials are multiscale, structurally inherited, and often directionally organized solids. Their response to heating and deformation reflects both intrinsic material properties and the geological architecture in which those properties are embedded. The next step is therefore to consider representative rock types whose contrasting origins and textures lead to contrasting thermoelastic behavior.